\section{Purity as Metaphysical Value}

\begin{quotex}
This is the introduction from the article titled “Purity as a Metaphysical Value” by \textbf{Julius Evola}. It was published in \textit{Bilychnis}, the journal of the Baptist Theological School of Rome. It is dated June, 1925, volume XXV. 

\end{quotex} 

There are two distinct ways that man can try to transcend the contingency and misery of the mundane world in relationship with the divinity. In the first, we presume that God is distinct from man, in a way that the relationship can only be extrinsic as in faith, prayer, devotion, or the observance of determinate moral principles to which one recognizes a higher validity. In the other, we postulate instead an ideal continuity between man and God, in a way that the relationship has the meaning of a real identification, of a joining of man to God, not with words, thoughts, or feelings, but rather by making oneself the God himself. This is the mystical and esoteric way, in contrast to the way of devotional religion. In the latter, therefore, the human state of existence is accepted and kept also close to faith in existence and a higher law; in the other, one would instead desire to truly transform without residue such a state, made of death and darkness, into the glory of a divine life.

In relationship to this difference, there are two absolutely distinct modes of realization of the concept of “purification”. Here we will be concerned only with the esoteric doctrine of purification, a rather impressive doctrine, about which however current culture knows only little or nothing. We will not dwell on the question of sources: it they are essentially connected to the Orient, and in particular to the tantric school and to the magical and alchemical offshoots of Taoism and Mahayana Buddhism, allusions to Greek mystery wisdom and pre-Socratic philosophy, elements of neoplatonism and a certain Christian mysticism—beyond formulations by no means different from what is present in the Kabbalistic, Hermetic, and Rosicrucian traditions.

We will not dwell on that both for reasons of space as well as the great difficulty of justifying with fresh cultural references that insofar as it is understood, to that extent a given interior attitude allows one to read between the lines, essentially for this reason: what matters is to present the doctrine in its logico-metaphysical essence, so that it constitutes something standing for itself, independent of beliefs, opinions, and various empirical elements that may have incorporated it.

We will immediately say that in the esoteric field the concept of purification has absolutely nothing moralistic about it. Instead, it is a question of a metaphysical value that, in its naked positivity, brings us back to the very literal meaning of the word. That which is not simply itself and from itself (kath'auto) must be called \emph{impure}, in general, i.e., it is contaminated by something else. Wherever, there is another, there is impurity: the other adulterates the being and renders it impure.

A passage from Eckhart can make that clear:

\begin{quotex}
Let's suppose that I take a burning coal and put it in my hand. If I said that it is the coal that burns my hand, I would be wrong. If I must truly indicate what burns me, I must instead say: it is the “nothing” that burns me! Since the coal has something in itself that my hand does not have! Precisely this “nothing” burns me. If my hand, instead, had in itself everything that the coal is and produces, it would entirely possess the nature of fire: and then when I also took all the fire that has ever burned and put it in my hand, it could no longer harm me 

\end{quotex}
That is, the being that is sufficient in the totality of life, would not have another against itself. Enclosed in an intangible unity, he would rest there and would indulge in it, loving himself alone and creating through this solitary love everything that creates. The point of the sufficient fails—then the unity is adulterated, a “nothing” is acknowledged so that, against the identity, the lauton, the other rises, the eteron. This other is not therefore anything real in itself; it is simply the reflection and the symbol of that deficiency that is produced in the being. Its substance being absolutely negative—not living for itself, but through the corruption of the perfect life—it is contingent in itself, that is, it subsists only in as much as, and as long as, the state of privation and imperfection remains in the being or in the I.

From this, the meaning of purifying: to purify means to bring life to the level of sufficient existence, of possession, of autarky, burning the dark privation of which it, at the point of finite and individual existence, is soaked and suffers from violence. From that, the concept of the impure act proceeds, which is clearly borrowed from the Aristotelian concept of the imperfect act: the impure is therefore the act of that power that does not reach actuality from itself, but that are needy of the complicity of the “other”. Such is, in order to immediately provide an example, the visual act, since in it the power of seeing is not sufficient in itself, does not produce vision from itself, but it needs a connection with a sensible object. What we will say below will clarify these points that perhaps now are expressed a little too abstractly.

The point on which, however, it is fundamentally necessary to focus attention, if one wants to understand the true meaning of the need of purification, is this: that the impure (or imperfect) act resolves the deficiency of the agent only apparently—it in reality reconfirms it. Man, for example, is thirsty: until he drinks, he will continue to be thirsty, since by drinking he confirms the state of one who does not have in himself his own life (to autarkes), but instead asks it from another. But this other—water and the rest—is only the symbol of his deficiency, and inasmuch as he feeds on it, and it demands his life, he in truth feeds only his own privation and remains in it avoiding that pure act, that eternal water through which every thirst, just as every other privation, would be forever expunged. Hence Christ says:

\begin{quotex}
Jesus answered, and said to her: Whosoever drinketh of this water, shall thirst again; but he that shall drink of the water that I will give him, shall not thirst for ever. But the water that I will give him, shall become in him a fountain of water, springing up into life everlasting. (John 4:13-14) 

\end{quotex}
Therefore every impure act is a flight from the perfect life; through it the individual does not endure but surrenders. Innumerable dependencies then condition him and confine him to the realm of contingency and death, of the Platonic “thing that is and is not”, and so he goes, passive to himself, dispersed in procreation and corruptions, in the eternal wheel of periodic, always different, rebirth and forever equal in their misery.

Now that the general meaning of the doctrine of the pure and the impure has been pointed out, we can move on to specify it in relation to particular problems, and effectively to say what the purification of the mind, the will, the word, the breath, and the generative act means in esoterism.



\flrightit{Posted on 2013-11-05 by Aeneas }
